\pg{352} \emph{simples}. Logo o ente, que é \emph{necessario} e \emph{infinito}, é absolutamente simples \footnote{A verdade d'estas proposições é tambem manifesta pela simples inspecção dos termos. S. Thomaz prova que o ente io é simples, porque o ente \emph{composto}, emquanto \emph{composto}, pode estar sujeito a dissolução, e por isso pode deixar de existir. (\emph{C}. Gent., I, 18).}.

77. O ente, que é composto, é contingente e finito; o ente, que é contingente e finito, é composto. :

a) \emph{O ente}, que é composto, é \emph{contingente} e \emph{finito}. --- .O ente, que é \emph{composto}, não é puro acto, mas resulta de potencia e de acto. Ora o ente, que resulta de potencia e de acto, --- é \emph{contingente}, porque não existe pela exigencia da sua essencia, mas é indifferente quanto á existencia, pode ser e deixar de ser, --- e é \emph{finito}, porque é composto de partes, essencialmente finitas. Logo o ente, que é \emph{composto}, \emph{é contingente} e \emph{finito}.

b) \emph{O ente}, que é \emph{contingente} e \emph{finito}, \emph{é composto}. --- O ente, que é contingente e \emph{finito}, possue a existencia realmente distincta da; essencia e resulta-de varias partes. Ora o ente, que possue a existencia realmente distincta da essencia e resulta de varias partes, é \emph{composto}; porque consta de potencia e de acto. Logo o ente, que é \emph{contingente} e \emph{finito}, \emph{é composto} \footnote{O \emph{composto é contingente}; porque, como diz S. Thomaz, o \emph{composto}, emquanto tal, está sujeito a dissolução e por isso pode deixar de ser. --- Noutro logar o 8. Doutor escreve: « Omne id, in cujus substantia admiscetur potenlia, secundum id quod habet de potentia, potest non esse; quia quod potest esse, potest non esse» (\emph{C}. \emph{Gent}., I, 16).}. '

78. O ente composto depende de uma \textbf{causa extrinseca}. --- O ente \emph{composto} .resulta da união de varias partes. A união, por ser um effeito, exige uma causa. Esta causa não pode ser o proprio \emph{composto}; porque. aliás, este, sendo naturalmente posterior ás partes componentes, operaria antes de existir e seria a causa de si mesmo. Deve ser, pois, uma causa \emph{extrinseca}. Logo o ente \emph{composto} depende de uma causa \emph{ezxirinseca} \footnote{O ente \emph{composto} é naturalmente posterior ás partes, de que resulta e que são diversas entre si, e por isso é um effeito. Qual é a causa d' este effeito? -- São as proprias partes? Não; porque \emph{partes} diversas, --- emquanto \emph{paries}, estão todas em potencia com relação ao todo, --- e, emquanto \emph{diversas}, \emph{não podem} por si formar a união. --- É o proprio composto ? Tambem não;\pgn{353}“porque, então, o composto seria causa de si mesmo, e existivia antes de existir. Logo deve ser uma causa \emph{extrinseca}. Diz S. Thomaz: «Omne compositum causam habet, Quae enim secundum se diversa sunt, non conveniunt in aliquod unum, nisi per aliquam causam adunantem ipsa » (\emph{Sum}. \emph{Th}., P. 1, q. 3, a. 7).}.

